\documentclass{scrartcl}
\usepackage[utf8]{inputenc}
\usepackage{fancyhdr}
\usepackage{graphicx}
\usepackage{dirtytalk}
\usepackage[portrait, margin=1in]{geometry}
\usepackage{amsmath, amssymb}
\setlength{\parskip}{1em}
\setlength{\parindent}{0cm} 
\usepackage[mdthm,simplethm]{test}
\usepackage[utf8]{inputenc}
\usepackage{float}
\usepackage{physics}
\setlength{\parindent}{0pt}

\title{Programs as Data}
\author{albert ye}
\date{May 2022}

\begin{document}
\maketitle
\section{Expressions as Lists}
\say{Lisp: everything's a list} -chez

Scheme programs are made up of expressions, which can be 
primitive expressions or combinations. But we can think of 
combinations as lists of primitive expressions. 

So basically, everything in Lisp is a List!

We have \texttt{(list 'quotient 10 2)} evaluate to \texttt{(quotient 10 2)}
while \texttt{(eval (list 'quotient 10 2))} evaluates to the value of the
quoted function, which is 5.

The takeaway here is that we can write programs that write programs,
which is hella cool

\section{Quasiquotation}
Quotations that can be unquoted, for example: 

\begin{verbatim}
	(define b 4)
	; the ` is a backtick dont forget
	`(a, (+ b 1)) => '(a (unquote (+ b 1))) => (a 5)
\end{verbatim}

\end{document}